\practicechapter{第 4 章实践：Linux、工具链、调试器和学习工作流}

本章对应第一册第 4 章“Linux、工具链、调试器和学习工作流”。正文讲命令行、编译器、调试器、readelf、objdump、perf 和可复现实验。本章把这些工具固定成一个日常工作流：先构建，再测试，再用工具交叉验证，最后把证据写入报告。

\practicesection{一、对应原书问题：工具不是截图生成器}

初学者常见做法是运行一堆命令，把终端输出截图放进文档，然后写“如图所示”。这不够。工具输出必须回答具体问题：\cmd{readelf -h} 验证 ELF header，\cmd{readelf -S} 验证 section table，\cmd{readelf -s} 或 \cmd{objdump -t} 验证 symbol，\cmd{objdump -d} 验证反汇编，\cmd{ctest} 验证代码合同。

本章实践目标是建立一套可重复命令，而不是临时点几下。每个命令都要写入报告，必要时把输出保存到 \filepath{reports/} 或 \filepath{results/}，这样别人才能复跑。

\practicesection{二、从特例推到规律：一个工具输出不能独占真相}

特例：\code{cpu1ee\_cli} 和 \cmd{readelf} 都能报告 machine。如果两者一致，说明我们对该字段的解析更可信；如果不一致，不能立刻说系统工具错了，更可能是我们的解析器、输入文件或字段理解有问题。另一个特例是 \cmd{objdump -d} 能看到反汇编，但它不能替你证明 benchmark 可信。

由此推出规律：工具要按问题组合使用。结构问题用 readelf 和本卷解析器交叉验证；指令问题用 objdump；运行问题用 perf 或 benchmark；正确性问题用 GTest/GMock；构建问题用 CMake 和 ctest。工具越多不代表越严谨，能互相验证同一个事实才严谨。

\practicesection{三、实践任务：建立固定工作流}

本章要求你把下面命令保存成自己的实验记录模板：

\begin{lstlisting}[language=bash]
cmake -S books/cpu-volume-1-practice \
      -B books/cpu-volume-1-practice/build/executable-evidence-debug \
      -DCMAKE_BUILD_TYPE=Debug
cmake --build books/cpu-volume-1-practice/build/executable-evidence-debug
ctest --test-dir books/cpu-volume-1-practice/build/executable-evidence-debug \
      --output-on-failure

cpu1ee=books/cpu-volume-1-practice/build/executable-evidence-debug/labs/executable_evidence/cpu1ee_cli
$cpu1ee "$binary" > "reports/$(basename "$binary").report"
$cpu1ee "$binary" --sections-csv > "results/$(basename "$binary").sections.csv"
readelf -h "$binary" > "reports/$(basename "$binary").readelf-header.txt"
readelf -S "$binary" > "reports/$(basename "$binary").readelf-sections.txt"
objdump -d "$binary" > "reports/$(basename "$binary").disasm.txt"
\end{lstlisting}

然后做一次字段核对：从 \code{cpu1ee\_cli} 报告中取 machine 和 entry，从 readelf header 中找同一字段；从 section CSV 中取 \code{.text} 的 offset 和 size，从 \cmd{readelf -S} 中找对应行。若格式不同，报告里写字段名如何对应。

\practicesection{四、验收：工作流要能给别人复跑}

本章交付物是 \filepath{reports/ch04-linux-tool-workflow.md}。通过标准如下：

\begin{enumerate}
  \item 报告列出构建、测试、CLI、readelf、objdump 的完整命令。
  \item 至少有三个字段被两个独立工具交叉验证。
  \item 能说明 \cmd{ctest} 验证的是代码行为，不验证真实二进制分析报告是否写得正确。
  \item 若当前环境缺少 \cmd{perf} 或权限不足，报告写明限制，不把缺失数据伪装成结论。
  \item 所有输出文件放在 \filepath{reports/} 或 \filepath{results/}，命名能看出对应的输入二进制。
\end{enumerate}

这章的重点不是记住每个命令选项，而是形成一个严谨习惯：每次观察都有问题、有命令、有输出、有对照、有边界。
